\chapter{Anhang A: Befehlsübersicht}

\begin{tabularx}{\textwidth}{@{}lX@{}}
\toprule
\textbf{Befehl} & \textbf{Bedeutung} \\
\midrule
\texttt{pwd}        & Aktuelles Verzeichnis anzeigen \\
\texttt{ls}         & Dateien und Ordner anzeigen \\
\texttt{ls -la}     & Mit Berechtigungen und versteckten Dateien \\
\texttt{cd}         & Verzeichnis wechseln \\
\texttt{mkdir}      & Verzeichnis erstellen \\
\texttt{touch}      & Datei erstellen \\
\texttt{cp}         & Datei oder Ordner kopieren \\
\texttt{mv}         & Datei oder Ordner verschieben / umbenennen \\
\texttt{rm}         & Datei löschen \\
\texttt{rm -r}      & Verzeichnis rekursiv löschen \\
\texttt{cat}        & Dateiinhalt ausgeben \\
\texttt{less}       & Datei seitenweise anzeigen \\
\texttt{grep}       & Text in Dateien suchen \\
\texttt{ps aux}     & Alle laufenden Prozesse anzeigen \\
\texttt{top / htop} & Prozesse interaktiv beobachten \\
\texttt{kill}       & Prozess beenden \\
\texttt{chmod}      & Dateiberechtigungen ändern \\
\texttt{sudo}       & Befehl mit Root-Rechten ausführen \\
\texttt{apt install}& Paket installieren (Ubuntu/Debian) \\
\midrule
\texttt{git init}   & Git-Repository initialisieren \\
\texttt{git status} & Zustand des Repositories anzeigen \\
\texttt{git add}    & Dateien zur Versionierung vormerken \\
\texttt{git commit} & Änderungen festhalten \\
\texttt{git log}    & Commit-Historie anzeigen \\
\texttt{git diff}   & Änderungen anzeigen \\
\texttt{git branch} & Branches anzeigen / erstellen \\
\texttt{git merge}  & Branch zusammenführen \\
\midrule
\texttt{docker ps}           & Laufende Container anzeigen \\
\texttt{docker images}       & Lokale Images anzeigen \\
\texttt{docker run}          & Container starten \\
\texttt{docker exec -it}     & In laufenden Container einsteigen \\
\texttt{docker build}        & Image aus Dockerfile bauen \\
\texttt{docker-compose up}   & Stack starten \\
\texttt{docker-compose down} & Stack stoppen \\
\bottomrule
\end{tabularx}

\chapter{Anhang B: Kapitel-Template}

Dieses Template beschreibt die Struktur jedes Kapitels in Band 1.
Es kann direkt als Ausgangspunkt für neue Kapitel verwendet werden.

\begin{lstlisting}[style=text]
\chapter{Wir wollen <konkrete Aufgabe>}

\kapitelziel{ ... }
\kapitelstruktur

\section{Situation}
...

\begin{aufgabeBox} ... \end{aufgabeBox}

\section{Erster Lösungsversuch}
...

\begin{problemBox} ... \end{problemBox}

\section{Was wir dafür lernen müssen}
\begin{wissenBox} ... \end{wissenBox}

\section{Der pragmatische Weg}
...

\section{Schritt-für-Schritt-Lösung}
\begin{loesungBox} ... \end{loesungBox}

\section{Ergebnis prüfen}
...

\begin{merksatzBox} ... \end{merksatzBox}
\begin{fehlerBox} ... \end{fehlerBox}
\begin{uebungBox} ... \end{uebungBox}
\begin{quizBox} ... \end{quizBox}
\begin{haubeBox} ... \end{haubeBox}
\begin{robotikBox} ... \end{robotikBox}
\end{lstlisting}
